% ============================================================
% 纯答案册·课时03-方向向量与法向量 —— 工具/答案抽册器.py 抽册生成件（勿手改；第三档＝纯答案单独成册）
% 源件（只读）：C:/提示词/工作区/M3-第2章量产0913/成卷/导学件/课时03-方向向量与法向量/main.tex（md5 63ea3670ea2d7f4ce61134aa7c88b1c1）
%% 口径：题面不重印；逐题＝题号(印面号同正册)＋[答案]＋[详解]，值/详解照源件逐字
%       （钉值门硬断言：册值≡正册件内值≡值台账值tex，逐字全等）。册序＝正册装配序＝印面号 1..21 升序（同序同号）；键序断言基准＝题面库冻结 manifest canonical 键序。
%% 三档导出：整体 main-true／纯题 main-false（在产）｜本册＝纯答案册（本件）
%% 双档：ansbook-true.tex＝详解印本；ansbook-false.tex＝纯值速查本（详解不印）
% 编译：xelatex <壳> ×2，cwd＝本目录；sty＝qp-m3.sty 本地挂载（md5 7c3930362be8a0a2bdf21bbf8ac16573，锁校验）
% ============================================================
\documentclass[fontset=none]{ctexart}
\usepackage{qp-m3}
% —— 印面逐字符号路由补钉（承源件；− 走 TNR、▱ NSC 子块；禁改 sty 保 md5） ——
\xeCJKDeclareCharClass{Default}{"2212}%
\setCJKfamilyfont{nscblk}[Path=C:/提示词/工作区/字替对照-0909/variantF/fonts/,BoldFont=NSC-w600.ttf]{NSC-w500.ttf}%
\xeCJKDeclareSubCJKBlock{nscblk}{"25B1}%
\newcommand{\pxparallelogram}{{\CJKfamily{nscblk}▱}}%
% —— 断行微调（承母版实测档，三零门承重；\emergencystretch 不另设） ——
\xeCJKsetup{CJKglue={\hskip 0.10em plus 0.08em minus 0.05em}}%
% —— 纯答案册全线括线（长详解可跨栏断，承练习件全线括线口径；灰底盒不可跨栏断） ——
\ansblockgrayfalse
% —— 双档开关（false 档＝纯值速查：答案印、详解不印；件面开关，不动 sty 本体） ——
\newif\ifansbookdetail
\ifdefined\ansbookdetail
  \ifnum\ansbookdetail=0 \ansbookdetailfalse\else\ansbookdetailtrue\fi
\else
  \ansbookdetailtrue
\fi
% —— 页脚件名（承源件章名；件名＝<件型>答案册） ——
\renewcommand{\qpzhangming}{第二章\quad 平面解析几何}%
\renewcommand{\qpjianming}{导学件答案册}%

\begin{document}
\zhangtitle{第二章\quad 平面解析几何}
\xjkeshi{2.2.1 直线的方向向量与法向量（第2课时）}
\par\glueguard{1}\addvspace{2pt}\noindent\biaoqian{【纯答案册】}{\fangsong 题面不重印\quad 印面号与正册同号同序}\par\addvspace{2pt}

\begin{multicols}{2}
\emergencystretch=1em

\begin{ansblock}[2章-练-课时03-E3]
% ans:2章-练-课时03-E3
\ansitem{1}{方向向量如(1,−1)（答案不唯一）；k＝−1；θ＝135°}
\ifansbookdetail
\ansnote{详解}{\(\overrightarrow{AB}=(-5+2,\ 3-0)=(-3,3)\)是直线\(l\)的一个方向向量（与之共线的非零向量如\((1,-1)\)也可）．斜率\(k=(3-0)/(-5-(-2))=3/(-3)=-1\)（也可由方向向量\((1,-1)\)得\(k=-1/1=-1\)）．由\(\tan\theta=-1\)且\(\theta\in[0^{\circ},180^{\circ})\)，得\(\theta=135^{\circ}\)．}
\fi
\end{ansblock}

\begin{ansblock}[2章-练-课时03-E1]
% ans:2章-练-课时03-E1
\ansitem{2}{2；4/3}
\ifansbookdetail
\ansnote{详解}{由方向向量\((2,4)\)得直线\(l\)的斜率\(k=4/2=2\)．由两点斜率公式\(k=(m-(-2))/(3-m)=(m+2)/(3-m)\)（\(m\neq 3\)），令\((m+2)/(3-m)=2\)，得\(m+2=6-2m\)，\(3m=4\)，\(m=4/3\)．故斜率为\(2\)，\(m=4/3\)．}
\fi
\end{ansblock}

\begin{ansblock}[2章-练-课时03-E2]
% ans:2章-练-课时03-E2
\ansitem{3}{(1) 1或2；(2) 4/3}
\ifansbookdetail
\ansnote{详解}{(1) \(k_{AC}=[4-(-m-3)]/(-1-m)=(m+7)/(-m-1)\)（\(m\neq -1\)），\(k_{BC}=(4-m+1)/(-1-2)=(5-m)/(-3)\)．由\(k_{AC}=3k_{BC}\)：\((m+7)/(-m-1)=3\cdot(5-m)/(-3)=m-5\)，去分母得\(m+7=(m-5)(-m-1)=-m^{2}+4m+5\)，即\(m^{2}-3m+2=0\)，解得\(m=1\)或\(m=2\)，经检验均使两斜率存在，符合题意．\par\noindent (2) 设\(l_1\)的倾斜角为\(\alpha\)，则\(l_2\)的倾斜角为\(2\alpha\)．由方向向量\(\overrightarrow{n}=(2,1)\)得\(\tan\alpha=1/2\)，故\(l_2\)的斜率\(\tan 2\alpha=2\tan\alpha/(1-\tan^{2}\alpha)=1/(1-1/4)=4/3\)．}
\fi
\end{ansblock}

\begin{ansblock}[2章-练-课时03-E4]
% ans:2章-练-课时03-E4
\ansitem{4}{A、B、C不共线；A、B、D共线．}
\ifansbookdetail
\ansnote{详解}{\(\overrightarrow{AB}=(0+1,\ -1+3)=(1,2)\)，\(\overrightarrow{AC}=(1+1,\ 2+3)=(2,5)\)．因为\(1\times 5-2\times 2=1\neq 0\)，\(\overrightarrow{AB}\)与\(\overrightarrow{AC}\)不共线，且两向量有公共起点\(A\)，所以\(A\)、\(B\)、\(C\)不共线．\(\overrightarrow{AD}=(3+1,\ 5+3)=(4,8)=4\overrightarrow{AB}\)，\(\overrightarrow{AB}\)与\(\overrightarrow{AD}\)共线且有公共起点\(A\)，所以\(A\)、\(B\)、\(D\)共线．}
\fi
\end{ansblock}

\begin{ansblock}[2章-练-课时03-E5]
% ans:2章-练-课时03-E5
\ansitem{5}{真命题．证明见详解．}
\ifansbookdetail
\ansnote{详解}{必要性：若\(A\)、\(B\)、\(C\)共线，则直线上的向量\(\overrightarrow{AB}\)、\(\overrightarrow{BC}\)及与直线平行的向量都是该直线的方向向量，故\(\overrightarrow{AB}\)与\(\overrightarrow{BC}\)共线．\par\noindent 充分性：若\(\overrightarrow{AB}\)与\(\overrightarrow{BC}\)共线，则直线\(AB\)与直线\(BC\)的方向相同或相反；又两直线过公共点\(B\)，由过一点且方向确定的直线唯一，直线\(AB\)与直线\(BC\)是同一条直线，故\(A\)、\(B\)、\(C\)三点共线．综上命题为真．}
\fi
\end{ansblock}

\begin{ansblock}[2章-练-课时03-E12]
% ans:2章-练-课时03-E12
\ansitem{6}{B}
\ifansbookdetail
\ansnote{详解}{\(\overrightarrow{AB}=(3-1,\ 6-2)=(2,4)=2(1,2)\)，故\((1,2)\)为\(l\)的一个方向向量，选B．（辨析：A 中\((2,4)=k(2,1)\)无解，与\(\overrightarrow{AB}\)不成比例；C 中\((-2,1)\cdot(2,4)=-4+4=0\)，是\(l\)的法向量方向而非方向向量；D 中\((1,-2)\)与\(\overrightarrow{AB}\)不成比例，均排除．）}
\fi
\end{ansblock}

\begin{ansblock}[2章-练-课时03-E15]
% ans:2章-练-课时03-E15
\ansitem{7}{m=0}
\ifansbookdetail
\ansnote{详解}{直线\(Ax+By+C=0\)的方向向量可取\((B,\ -A)\)，故\(l\)的方向向量\(t=(2,\ -(m+1))\)．\(t\parallel v\)（横坐标已同为\(2\)）\(\Rightarrow -(m+1)=m-1\)\(\Rightarrow 2m=0\)\(\Rightarrow m=0\)．（验证：\(m=0\)时\(l\)为\(x+2y-6=0\)，方向向量\((2,-1)=v\)；又法向量\((1,2)\)，\(v\cdot n=2-2=0\)，两法一致．）}
\fi
\end{ansblock}

\begin{ansblock}[2章-练-课时03-E8]
% ans:2章-练-课时03-E8
\ansitem{8}{一定．理由见详解．}
\ifansbookdetail
\ansnote{详解}{设\((x_0,y_0)\)是\(l\)的一个方向向量．由方向向量非零知\((x_0,y_0)\neq(0,0)\)，从而\((-y_0,x_0)\neq(0,0)\)．又\((x_0,y_0)\cdot(-y_0,x_0)=-x_0y_0+y_0x_0=0\)，即\((-y_0,x_0)\)与\(l\)的方向向量垂直，于是以\((-y_0,x_0)\)为方向向量的有向线段所在直线与\(l\)垂直，由法向量的定义，\((-y_0,x_0)\)是\(l\)的一个法向量．反之，若\((-y_0,x_0)\)为\(l\)的方向向量，同理\((-y_0,x_0)\neq(0,0)\)即\((x_0,y_0)\neq(0,0)\)，且两向量数量积为\(0\)，故\((x_0,y_0)\)一定是\(l\)的法向量．所以结论成立．}
\fi
\end{ansblock}

\begin{ansblock}[2章-练-课时03-E11]
% ans:2章-练-课时03-E11
\ansitem{9}{证明见详解．}
\ifansbookdetail
\ansnote{详解}{\(\overrightarrow{AB}=(3-1,\ 3-4)=(2,-1)\)，故直线\(AB\)的一个方向向量为\(a=(2,-1)\)；\(\overrightarrow{BC}=(4-3,\ 5-3)=(1,2)\)．取\(v=(1,2)\)，验证\(v\)与\(a\)垂直：\(v\cdot a=1\times 2+2\times(-1)=0\)，且\(v\neq 0\)，所以\(v\)是直线\(AB\)的一个法向量．而\(\overrightarrow{BC}=v\)，即直线\(BC\)的方向向量就是直线\(AB\)的法向量，由法向量的定义，直线\(BC\)与直线\(AB\)垂直，故\(AB\perp BC\)．}
\fi
\end{ansblock}

\begin{ansblock}[2章-练-课时03-E14]
% ans:2章-练-课时03-E14
\ansitem{10}{(1) n₁=(3,−1)（任一非零倍数均可）　(2) a=2/3}
\ifansbookdetail
\ansnote{详解}{(1) 由一般式系数直接读出\(n_1=(3,-1)\)．(2) \(n_2=(a,2)\)．同一平面内每条直线的法向量与其方向向量互相垂直，将两直线绕任一点同旋\(90^{\circ}\)，方向向量恰转到各自法向量方向且夹角不变，故\(l_1\perp l_2\)\(\iff n_1\perp n_2\)\(\iff n_1\cdot n_2=0\)\(\iff 3a+(-1)\times 2=0\)\(\iff a=2/3\)．}
\fi
\end{ansblock}

\begin{ansblock}[2章-练-课时03-E13]
% ans:2章-练-课时03-E13
\ansitem{11}{2x−y−5=0}
\ifansbookdetail
\ansnote{详解}{法向量\((2,-1)\)\(\Rightarrow l\)具形式\(2x-y+C=0\)．代\(P(3,1)\)：\(6-1+C=0\)\(\Rightarrow C=-5\)．故\(2x-y-5=0\)．（验证：\(P\)代入\(6-1-5=0\)；\(l\)的方向向量\((1,2)\)，\(n\cdot(1,2)=2-2=0\)．）}
\fi
\end{ansblock}

\begin{ansblock}[2章-练-课时03-E6]
% ans:2章-练-课时03-E6
\ansitem{12}{(1) y＝−2x−4（即2x＋y＋4＝0）；(2) 当u≠0时为y−y₀＝(v/u)(x−x₀)；当u＝0时为x＝x₀．}
\ifansbookdetail
\ansnote{详解}{(1) 方向向量\(a=(1,-2)\)的横坐标\(1\neq 0\)，故斜率\(k=-2/1=-2\)，由点斜式得\(y-(-2)=-2[x-(-1)]\)，即\(y=-2x-4\)．\par\noindent (2) 设\(Q(x,y)\)是直线\(l\)上任意一点，则\(\overrightarrow{PQ}=(x-x_0,\ y-y_0)\)与方向向量\(a=(u,v)\)共线，得\(v(x-x_0)-u(y-y_0)=0\)．当\(u\neq 0\)时化为\(y-y_0=(v/u)(x-x_0)\)；当\(u=0\)时（此时\(v\neq 0\)）得\(x-x_0=0\)，即\(x=x_0\)（斜率不存在的竖直线）．}
\fi
\end{ansblock}

\begin{ansblock}[2章-练-课时03-E7]
% ans:2章-练-课时03-E7
\ansitem{13}{(1) y＝3x；(2) u(x−x₀)＋v(y−y₀)＝0，即ux＋vy−ux₀−vy₀＝0．}
\ifansbookdetail
\ansnote{详解}{(1) 设\(Q(x,y)\)是直线\(l\)上任意一点，则\(\overrightarrow{PQ}=(x-1,\ y-3)\)与法向量\(a=(-3,1)\)垂直，\(a\cdot\overrightarrow{PQ}=0\)，即\(-3(x-1)+1\cdot(y-3)=0\)，化简得\(y=3x\)．\par\noindent (2) 同法：\(a\cdot\overrightarrow{PQ}=u(x-x_0)+v(y-y_0)=0\)（\(u\)、\(v\)不全为\(0\)，因法向量是非零向量）．这就是直线的点法式方程；当\(v\neq 0\)时可化为斜截形\(y=-(u/v)x+(ux_0+vy_0)/v\)，斜率\(k=-u/v\)（与"方向向量\((v,-u)\)读斜率\(-u/v\)"一致）．}
\fi
\end{ansblock}

\begin{ansblock}[2章-练-课时03-E9]
% ans:2章-练-课时03-E9
\ansitem{14}{(1) y＝2x−11；(2) y＝−(4/3)x＋5（即4x＋3y−15＝0）}
\ifansbookdetail
\ansnote{详解}{(1) \(n=(1,2)\)横坐标不为\(0\)，\(k=2/1=2\)，由点斜式\(y+5=2(x-3)\)，即\(y=2x-11\)．\par\noindent (2) \(n=(3,-4)\)横坐标不为\(0\)，\(k=-4/3\)，由点斜式\(y-5=-(4/3)(x-0)\)，即\(y=-(4/3)x+5\)．}
\fi
\end{ansblock}

\begin{ansblock}[2章-练-课时03-E10]
% ans:2章-练-课时03-E10
\ansitem{15}{(1) 3x−4y＋5＝0；(2) 3x＋4y−5＝0}
\ifansbookdetail
\ansnote{详解}{(1) 设\(Q(x,y)\)在\(l\)上，\(v\cdot\overrightarrow{PQ}=3(x-1)-4(y-2)=0\)，化简得\(3x-4y+5=0\)．\par\noindent (2) 同法：\(3(x+1)+4(y-2)=0\)，化简得\(3x+4y-5=0\)．}
\fi
\end{ansblock}

\begin{ansblock}[2章-练-课时03-E16]
% ans:2章-练-课时03-E16
\ansitem{16}{−√3/3}
\ifansbookdetail
\ansnote{详解}{斜率\(k=\sqrt{3}/1=\sqrt{3}\)，\(l\)：\(y=\sqrt{3}x+1\)．令\(y=0\)：\(x=-1/\sqrt{3}=-\sqrt{3}/3\)．（验证：点\((-\sqrt{3}/3,\ 0)\)代入\(\sqrt{3}\times(-\sqrt{3}/3)+1=-1+1=0\)，斜率与\(v\)一致．）}
\fi
\end{ansblock}

\begin{ansblock}[2章-导-课时03-G1]
% ans:2章-导-课时03-G1
\ansitem{17}{D}
\ifansbookdetail
\ansnote{详解}{直线\(l\)的一个方向向量为\(\overrightarrow{PQ}=(-2-1,\ -2-2)=(-3,-4)\)，\(|\overrightarrow{PQ}|=\sqrt{(-3)^{2}+(-4)^{2}}=5\)，故单位方向向量为\(\pm\overrightarrow{PQ}/|\overrightarrow{PQ}|=\pm(1/5)(-3,-4)=\pm(3/5,4/5)\)（方向相同或相反的两个单位向量都是方向向量）．C 写成"\(\pm\)"作用于每个分量的四种组合，其中\((3/5,-4/5)\)不与\(\overrightarrow{PQ}\)共线，故C错．选D．}
\fi
\end{ansblock}

\begin{ansblock}[2章-导-课时03-G2]
% ans:2章-导-课时03-G2
\ansitem{18}{B}
\ifansbookdetail
\ansnote{详解}{\(\overrightarrow{PQ}=(-m-3,\ 2-2m)\)．两直线平行\(\iff\)方向向量共线（此处即向量共线，注意\(P\)、\(Q\)不重合需\(2m\neq 2\)即\(m\neq 1\)），由共线条件\((-m-3)\times 5-(2-2m)\times(-5)=0\)，得\(-5m-15+10-10m=0\)，\(-15m-5=0\)，\(m=-1/3\)．经检验\(m=-1/3\)符合题意．选B．}
\fi
\end{ansblock}

\begin{ansblock}[2章-导-课时03-G3]
% ans:2章-导-课时03-G3
\ansitem{19}{A}
\ifansbookdetail
\ansnote{详解}{法向量\(v=(2,-4)\)垂直于直线\(l\)，故\(l\)的一个方向向量为与\(v\)垂直的非零向量\(a=(4,2)\)（因\(v\cdot a=2\times 4+(-4)\times 2=0\)），其横坐标不为\(0\)，斜率\(k=2/4=1/2\)．故选A．（另解：法向量为\((A,B)\)的直线可写成\(Ax+By+C=0\)形，\(2x-4y+C=0\)即\(y=x/2+C/4\)，\(k=1/2\)．）}
\fi
\end{ansblock}

\begin{ansblock}[2章-导-课时03-G4]
% ans:2章-导-课时03-G4
\ansitem{20}{(3,4)（答案不唯一）；5}
\ifansbookdetail
\ansnote{详解}{直线\(AB\)的一个方向向量为\(\overrightarrow{AB}=(3-0,\ 0+4)=(3,4)\)（与之平行的任何非零向量均可）；线段\(AB\)的长度为\(|AB|=\sqrt{3^{2}+4^{2}}=5\)．}
\fi
\end{ansblock}

\begin{ansblock}[2章-导-课时03-G5]
% ans:2章-导-课时03-G5
\ansitem{21}{[0,3]}
\ifansbookdetail
\ansnote{详解}{方向向量的横坐标为\(1\neq 0\)，故直线\(AB\)的斜率存在，设为\(k\)．由方向向量\((1,\sin\alpha)\)知\(k=\sin\alpha/1=\sin\alpha\in[-1,1]\)；又由两点坐标\(k=(2m-1-2)/(-2-1)=(3-2m)/3\)．所以\(-1\leqslant(3-2m)/3\leqslant 1\)，解得\(0\leqslant m\leqslant 3\)．故\(m\)的取值范围为\([0,3]\)．}
\fi
\end{ansblock}

% ---- 尾块（\tailfill[30mm] 定高参——钉档 §三.3：无参在平衡栏呈「内容尾随框」，贴栏底须定高参；实测无参致末页平衡 Overfull\vbox 12.99pt，定高 30mm 三零） ----
\tailfill[30mm]

\end{multicols}
\end{document}
